\documentclass[10pt,letterpaper,twocolumn]{article}

% ==================== PAQUETES ====================
\usepackage[utf8]{inputenc}
\usepackage[T1]{fontenc}
\usepackage[spanish]{babel}
\usepackage[margin=1cm,bottom=2cm,columnsep=0.7cm]{geometry}
\usepackage{graphicx}
\usepackage{booktabs}
\usepackage{array}
\usepackage{enumitem}
\usepackage{hyperref}
\usepackage{csquotes}
\usepackage{titlesec}
\usepackage{abstract}
\usepackage{times}    
\usepackage{caption}

% --- estilo de títulos de sección, más compacto ---
\titleformat{\section}
  {\normalfont\large\bfseries}{\thesection.}{0.5em}{}
\titleformat{\subsection}
  {\normalfont\normalsize\bfseries\itshape}{\thesubsection.}{0.5em}{}

% --- formato del título y abstract ---
\renewcommand{\abstractnamefont}{\normalfont\bfseries}
\renewcommand{\abstracttextfont}{\normalfont\small\itshape}
\setlength{\absleftindent}{0.5cm}
\setlength{\absrightindent}{0.5cm}

% --- espaciado de párrafos ---
\setlength{\parskip}{0.3em}
\setlength{\parindent}{1em}
\setlength{\columnsep}{0.7cm}

% --- captions más pequeños ---
\captionsetup{font=small,labelfont=bf}

\hypersetup{
  colorlinks=true,
  linkcolor=black,
  urlcolor=blue,
  citecolor=black
}

% ==================== DATOS DEL TRABAJO ====================
\title{\textbf{\Large Historia de las Computadoras: \\
de la Máquina Analítica a la Quinta Generación}}

\author{Bryan Ezequiel Violante Arriaga}

\date{\small\today}

\begin{document}

% ==================== ENCABEZADO A UNA COLUMNA ====================
\twocolumn[
\begin{@twocolumnfalse}
  \maketitle
  \begin{abstract}
    \noindent
    El presente trabajo (correspondiente a la actividad 1 del curso 
    propedéutico de Programación y Estructuras de Datos para la MC. en Ciencias 
    de la Computación) expone la evolución histórica de las computadoras,
    desde los primeros antecedentes mecánicos del siglo~XIX hasta los
    sistemas de quinta generación. La revisión se organiza siguiendo la
    convención de \emph{generaciones}, asociadas a saltos tecnológicos
    sucesivos: tubos de vacío, transistores, circuitos integrados,
    microprocesadores y, finalmente, sistemas paralelos orientados a
    inteligencia artificial, mencionando de manera breve la Computación Cuántica. 
    \vspace{0.5em}
  \end{abstract}
  
\end{@twocolumnfalse}
]

% ============================================================
\section{Introducción}
% ============================================================
Cualquier persona interesada en la informática, más generalmente en las ciencias
de la computación, debería conocer la historia de las computadoras. No solo por 
una cuestión de cultura general, sino porque entender cómo llegamos a donde 
estamos hoy es clave para comprender los conceptos y tecnologías que usamos a 
diario, y para anticipar hacia dónde podríamos ir en el futuro.

La fuente principal de esta investigación fue el libro de Andrew~S. Tanenbaum
\emph{Sistemas Operativos Modernos}, que dedica una sección a la historia de las 
computadoras. Para complementar esta fuente, se consultaron también diversas 
fuentes académicas en línea.  \\
La división en generaciones es una convención que se usa a menudo para organizar 
la historia, y básicamente se asocia a saltos tecnológicos sucesivos. 
Sin embargo, como advierte el propio Tanenbaum, esta división es ``un poco burda'' 
y las fases se traslapan, con falsos inicios y callejones sin salida. 

Antes de comenzar definiremos brevemente el concepto de computadora que, al ser 
un elemento fundamental en nuestra vida diaria, muchas veces damos por sentado.

\subsection{¿Qué es una computadora?}
Una computadora, computador, u ordenador es una máquina programable para el 
procesamiento de información. En otras palabras, una computadora es un 
dispositivo que puede recibir datos de entrada, procesarlos (mediante 
algoritmos) y generar resultados de salida. \\
En la vida cotidiana, cada vez que se hace referencia al término ``computadora'' 
se hace referencia de manera implícita a las computadoras digitales pero no es el
único tipo de computadora existente.\\
Las computadoras clasifican principalmente de acuerdo al 
principio de operación en Analógicas y Digitales.

\subsection{Computadoras Analógicas}
La computadora analógica es un dispositivo electrónico o hidráulico diseñado 
para manipular la entrada de datos en términos de niveles de 
tensión o presiones hidráulicas, en lugar de hacerlo como datos numéricos. 
El dispositivo de cálculo analógico más sencillo es la regla de cálculo, que 
utiliza longitudes de escalas especialmente calibradas para facilitar 
la multiplicación, la división y otras funciones.

\subsection{Computadoras Digitales}
Como se mencionó, las computadoras son dispositivos capaces de recibir 
información, 
procesarla (mediante algoritmos) y generar resultados de salida. A diferencia de
las computadoras analógicas, que manipulan datos en forma continua, 
las computadoras digitales representan la información en forma discreta,
usualmente en formato binario (0s y 1s).


% ============================================================
\section{Antecedentes: antes de la era electrónica}
% ============================================================

Algunas fuentes consideran que la historia de las computadoras comienza con la
invención del ábaco, un dispositivo de cálculo mecánico utilizado desde la 
antigüedad para realizar operaciones aritméticas básicas por los griegos y los
romanos. \\
Otro antecedente importante es la \textbf{máquina de Pascal} o 
\textbf{Pascalina}, 
inventada por Blaise Pascal en 1642, que era capaz de realizar sumas y restas 
utilizando ruedas dentadas. \\ 
En este trabajo (como en la mayoría de los trabajos revisados) se considera que 
la historia de las computadoras si bien no comienza con la electrónica, 
comienza con el proyecto del matemático inglés 
\textbf{Charles Babbage} (1792--1871) considerado como algunos 
como el padre de la computación, en el siglo~XIX . \\ 
Charles Babbage diseñó la \textbf{máquina analítica}, la cual es 
considerada por la literatura como la primera computadora digital concebida en 
la historia. A pesar de que Babbage dedicó gran parte de sus recursos y tiempo 
a este proyecto, nunca logró que funcionara, debido a que la tecnología de su 
época no podía fabricar los elementos con la precisión que el diseño exigía. 
Estaba adelantado a su tiempo.

En la definición de computadora se mencionó vagamente el concepto de 
\textbf{algoritmo} que, en pocas palabras es una serie de pasos para resolver 
un problema. Relacionada con este poderoso concepto al igual que con la máquina 
analítica de Babbage, aparece otra figura fundamental: considerada la primera 
programadora de la historia, \textbf{Ada Lovelace}, fue contratada por Babbage 
para escribir los algoritmos que ejecutaría su máquina (razón por la cual se 
nombró posteriormente al lenguaje de programación \texttt{Ada}). 

Aunque la máquina analítica nunca operó, su importancia es enorme: sentó
las bases conceptuales de una computadora programable de propósito general
más de un siglo antes de que la tecnología pudiera construirla.

% ============================================================
\section{Primera generación (1945--1955): tubos de vacío}
% ============================================================

\subsection{Contexto y tecnología base}

La construcción de computadoras digitales avanzó muy poco
hasta la WW2 (Segunda Guerra Mundial), que impulsó esta actividad por
necesidades militares: cálculos balísticos, criptoanálisis y simulaciones
físicas.

La tecnología característica fue el \textbf{tubo de vacío} (o bulbo).
Estas máquinas eran enormes, ocupaban habitaciones completas, consumían
cantidades enormes de energía y generaban mucho calor.
Una computadora típica podía contener cerca de \textbf{20{\,}000 bulbos},
lo que tenía inconvenientes bastante evidentes como el fundido regular de los 
bulbos durante la ejecución (lo que detenía el cómputo en curso).

\subsection{Modelos representativos}

\begin{itemize}[leftmargin=*,itemsep=0.2em]
  \item \textbf{Atanasoff-Berry Computer (ABC).} Construida por John
    Atanasoff y Clifford Berry en la Universidad Estatal de Iowa.
    Considerada la primera computadora digital electrónica funcional;
    utilizaba alrededor de 300 tubos de vacío.
  \item \textbf{Z3.} Construida por Konrad Zuse en Berlín, basada en
    relevadores electromecánicos.
  \item \textbf{Colossus (1944).} Desarrollada en Bletchley Park
    (Inglaterra) por un equipo dedicado a descifrar códigos militares
    alemanes durante la guerra.
  \item \textbf{Mark~I.} Construida por Howard Aiken en la Universidad
    de Harvard.
  \item \textbf{ENIAC.} Desarrollada por John Mauchly y J.~Presper
    Eckert en la Universidad de Pennsylvania. Es probablemente la más
    conocida del periodo.
\end{itemize}

\subsection{Cómo se usaban}

A diferencia de los recursos que tenemos hoy en día, no existían sistemas 
operativos ni lenguajes de programación, ni siquiera el tan temido lenguaje 
ensamblador. La programación se hacía directamente en \textbf{lenguaje máquina},
 o conectando físicamente miles de cables a \textbf{tableros de conexiones} 
 (\emph{plugboards}) para definir el comportamiento de la máquina. \\
Los problemas que se resolvían eran principalmente cálculos numéricos simples: 
tablas de senos, cosenos y logaritmos.

A principios de la década de 1950, la introducción de las
\textbf{tarjetas perforadas} mejoró parcialmente este flujo, permitiendo
escribir programas en tarjetas en lugar de cablear plugboards.

\subsection{Aporte clave}

Demostraron que el cómputo electrónico de propósito general era posible
y sentaron la base de toda la industria que vendría después.

% ============================================================
\section{Segunda generación (1955--1965): transistores}
% ============================================================

\subsection{Contexto y tecnología base}

A mediados de los años~50 apareció una innovación que es considerada como el 
invento del siglo y que cambió el panorama por completo: el \textbf{transistor}.
A diferencia de los bulbos, los transistores eran pequeños, 
consumían poca energía, generaban poco calor y eran mucho más confiables. 
Esto permitió, por primera vez, fabricar y vender computadoras a clientes con 
la expectativa razonable de que funcionaran de forma estable durante periodos 
útiles de trabajo.

El núcleo de esta etapa fueron los transistores discretos, las memorias de
núcleos magnéticos y las cintas magnéticas como medio de almacenamiento
secundario.

\subsection{Modelos representativos}

A finales de los 50's, principios de los 60's, las computadoras seguían 
evolucionando, reduciendo su tamaño y aumentando su potencia. \\
Estas máquinas conocidas como \textbf{mainframes} estaban encerradas en cuartos 
especiales con aire acondicionado y eran operadas por personal especializado.
Debido a su alto costo y complejidad de operación, solo grandes empresas, 
universidades y el gobierno podían darse el lujo de tener una. \\
Se desarrollaron dos líneas de máquinas, una orientada a aplicaciones 
comerciales y otra orientada a aplicaciones científicas.

\begin{itemize}[leftmargin=*,itemsep=0.2em]
  \item \textbf{IBM~1401.} Máquina comercial orientada a caracteres,
    ampliamente usada por bancos y aseguradoras para tareas como ordenar
    cintas e imprimir reportes. No era buena para cálculo numérico, pero
    sí para entrada/salida.
  \item \textbf{IBM~7094.} Máquina científica orientada a palabras,
    utilizada para cálculos numéricos pesados (ecuaciones diferenciales
    en física e ingeniería). Programada principalmente en
    \textsc{Fortran} y en lenguaje ensamblador.
\end{itemize}

\subsection{Cómo se usaban}

Por primera vez existía una clara separación entre quienes diseñaban,
construían, operaban y programaban las máquinas. Para ejecutar un
trabajo, el programador escribía el programa en papel y después lo pasaba a
tarjetas perforadas, las cuales eran entregadas en una ventanilla y horas
después regresaba por los resultados impresos. \\
Este flujo dio origen al concepto de \textbf{procesamiento por lotes} 
(\emph{batch processing}), --- que, al menos en experiencia del autor, es 
ampliamente utilizado en el contexto de machine learning, cuando los datos de 
entrenamiento y prueba se procesan en lotes --- 
en el que varios trabajos se acumulaban en una cinta y se ejecutaban uno tras 
otro sin intervención humana.

\subsection{Aporte clave}

La transición del tubo al transistor hizo que las computadoras fueran
\textbf{confiables, comercializables y mucho más eficientes}, abriendo
el mercado más allá de los laboratorios de investigación.

% ============================================================
\section{Tercera generación (1965--1980): circuitos integrados}
% ============================================================

\subsection{Contexto y tecnología base}

A inicios de los años~60, los fabricantes mantenían dos líneas de
productos incompatibles: máquinas científicas (como la 7094) y máquinas
comerciales (como la 1401). Mantener ambas era caro, y los clientes que
crecían querían poder migrar a máquinas más potentes sin reescribir todo
su software. Este problema impulsó uno de los lanzamientos más
importantes de la historia. El \textbf{circuito integrado} (IC) a pequeña
escala, que permitió empaquetar varios transistores en un solo chip de
silicio. El resultado fue una mejora notable en la relación
precio/rendimiento respecto a las máquinas de transistores discretos.

\subsection{Modelos representativos}

\begin{itemize}[leftmargin=*,itemsep=0.2em]
  \item \textbf{IBM System/360 (1964).} La primera \emph{familia} de
    computadoras compatibles entre sí en software. Iba desde modelos
    pequeños hasta modelos más potentes que la 7094, todos con la
    misma arquitectura y conjunto de instrucciones. Fue un éxito inmediato e
    instauró la idea de \emph{familia compatible}, que el resto de la
    industria adoptó.
  \item \textbf{Sucesores de la línea 360.} Modelos 370, 4300, 3080,
    3090 y, más recientemente, la serie zSeries. Sus descendientes
    siguen en uso en grandes centros de cómputo (sistemas de
    reservaciones aéreas, bases de datos masivas, servidores web de
    alto volumen).
  \item \textbf{DEC PDP-1 (1961).} Pionera de las minicomputadoras. 
    Tenía solo 4K palabras de 18 bits, pero
    a un precio de 120{\,}000~dólares (menos del 5\,\% del costo de
    una 7094). Fue un éxito de ventas y dio origen a una nueva
    industria.
  \item \textbf{DEC PDP-11.} Culminación de la línea de
    minicomputadoras de DEC, ampliamente influyente en el mundo
    académico.
\end{itemize}

\subsection{Cómo se usaban}

Tres ideas clave de software acompañaron a este hardware:
\textbf{multiprogramación} (mantener varios trabajos en memoria al mismo
tiempo para que la CPU no quedara ociosa cuando uno esperaba E/S),
\textbf{spooling} (leer trabajos directamente al disco para que el
sistema operativo pudiera tomar el siguiente automáticamente) y
\textbf{tiempo compartido} o \emph{timesharing} (permitir que varios
usuarios trabajaran simultáneamente desde terminales conectadas a la
misma máquina).

\subsection{Aporte clave}

La idea de \textbf{familia compatible de computadoras} cambió la
industria para siempre, y las minicomputadoras democratizaron el acceso
al cómputo fuera de las grandes corporaciones.

% ============================================================
\section{Cuarta generación (1980--ca.~2010): computadoras personales}
% ============================================================

\subsection{Contexto y tecnología base}

El desarrollo de los \textbf{circuitos LSI} (\emph{Large Scale
Integration}), capaces de empaquetar miles de transistores en un
centímetro cuadrado de silicio, dio nacimiento al
\textbf{microprocesador}, que es un adelanto enorme de la microelectrónica y, con 
él, a la era de la computadora personal.

\subsection{Procesadores y modelos representativos}

\begin{itemize}[leftmargin=*,itemsep=0.2em]
  \item \textbf{Intel~8080 (1974).} Primera CPU de 8~bits de propósito
    general. Marcó el inicio de la microcomputación moderna.
  \item \textbf{IBM~PC (1981).} Basada inicialmente en el procesador
    Intel~8088. Su arquitectura abierta y la decisión de IBM de no
    monopolizar el sistema operativo generaron un ecosistema enorme de
    fabricantes \emph{compatibles con IBM~PC}.
  \item \textbf{IBM~PC/AT (1983).} Sucesora basada en el Intel~80286.
  \item \textbf{Línea Intel x86.} 8086, 80286, 80386, 80486 y la
    familia \textbf{Pentium} (I, II, III, 4) y posteriormente Core~2~Duo
    y sus sucesores.
  \item \textbf{Apple Lisa (1983) y Apple Macintosh (1984).} Pioneras en
    llevar la \textbf{interfaz gráfica de usuario (GUI)} al mercado de
    consumo. La Lisa fue un fracaso comercial por su precio elevado,
    pero la Macintosh fue un enorme éxito por su precio más accesible y
    su enfoque amigable con el usuario.
  \item \textbf{Estaciones de trabajo RISC.} Sun Microsystems y
    Hewlett-Packard llevaron al mercado profesional máquinas basadas
    en arquitecturas RISC de alto rendimiento, estándar en aplicaciones
    técnicas y científicas.
\end{itemize}

\subsection{Cómo se usaban}

Hubo dos cambios profundos en la relación humano-máquina. El primero:
la computadora dejó de ser un recurso compartido y carísimo, y se
convirtió en una herramienta personal. El segundo: la \textbf{interfaz
gráfica de usuario} (con ventanas, íconos, menús y ratón), inventada
conceptualmente por Doug~Engelbart en Stanford y desarrollada en los
laboratorios de Xerox~PARC, fue adoptada por Apple y posteriormente por
Microsoft, transformando a la computadora en un dispositivo accesible
para usuarios sin formación técnica.

\subsection{Aporte clave}

La \textbf{democratización masiva} del cómputo: por primera vez, tener
una computadora propia se volvió realista para individuos, escuelas y
pequeñas empresas.

% ============================================================
\section{Quinta generación: dos visiones complementarias}
% ============================================================

A diferencia de las generaciones anteriores, no existe un consenso único
sobre qué constituye la \textbf{quinta generación} de computadoras. En
la literatura conviven dos definiciones complementarias.

\subsection{Visión clásica: el proyecto japonés FGCS (1982--1992)}

En 1982, el gobierno japonés, a través del Ministerio de Industria y
Comercio Internacional (MITI), lanzó el proyecto \textbf{Fifth
Generation Computer Systems (FGCS)}. Sus metas eran ambiciosas:
construir computadoras capaces de procesar \emph{conocimiento} y no
solo datos numéricos; apoyarse en \textbf{procesamiento paralelo
masivo} en lugar del procesamiento secuencial tradicional; usar
lenguajes de \textbf{programación lógica},
como base de la inteligencia artificial; e implementar interfaces de
lenguaje natural y reconocimiento de voz e imagen.

El proyecto se considera un \emph{éxito parcial}: no cumplió todas sus
metas (la IA simbólica resultó más difícil de lo previsto y los
lenguajes lógicos no se impusieron como modelo dominante), pero sembró
ideas que más tarde fructificarían en otros contextos.

\subsection{Visión moderna: paralelismo, móviles, nube e IA}

Una segunda definición, más actual, sitúa a la quinta generación en el
periodo que va desde la década de 1990 hasta el presente. En cuanto al
hardware, la tecnología base evolucionó del LSI/VLSI al
\textbf{ULSI} (\emph{Ultra Large Scale Integration}), que permite
empaquetar millones e incluso miles de millones de transistores en un
solo chip. Sobre esta base se construyen \textbf{procesadores
multinúcleo} y el cómputo paralelo se vuelve cotidiano en dispositivos
de uso diario.

Otros rasgos característicos son: \textbf{Internet} como
infraestructura central y el \textbf{cómputo en la nube} (\emph{cloud
computing}) como modelo dominante de servicios; el \textbf{cómputo
móvil} con \emph{smartphones} y tablets como computadoras personales
primarias para gran parte de la población mundial; y, sobre todo, la
consolidación de la \textbf{inteligencia artificial} como aplicación
central del cómputo.

La IA de esta generación se manifiesta en varias ramas concretas:
\textbf{sistemas expertos}, \textbf{redes neuronales} y aprendizaje
automático, \textbf{robótica}, \textbf{procesamiento de lenguaje
natural} y reconocimiento de voz e imagen. En la última década, los
modelos de aprendizaje profundo y los modelos de lenguaje a gran escala
han ampliado significativamente lo que estas máquinas pueden hacer.

\subsection{¿Hacia una sexta generación? La computación cuántica}

En la visión del autor se considera que la \textbf{computación cuántica}, junto
con otras tecnologías emergentes como la computación neuromórfica,
representa el inicio de una posible \emph{sexta generación} de
computadoras. La razón es que, a diferencia de todas las generaciones
anteriores (que son variaciones tecnológicas del mismo paradigma de
cómputo clásico), la computación cuántica representa un cambio de
paradigma físico fundamental.

Mientras una computadora clásica usa \textbf{bits} que valen 0 o 1, una
computadora cuántica usa \textbf{qubits} que pueden estar en una
superposición de ambos estados al mismo tiempo, gracias a fenómenos de
la mecánica cuántica como la superposición y el entrelazamiento. Esto
no implica que las computadoras cuánticas sean ``más rápidas'' en
general, sino que son potencialmente mucho más eficientes para una
clase específica de problemas.

Algunos hitos relevantes en esta área son:

\begin{itemize}[leftmargin=*,itemsep=0.2em]
  \item \textbf{1980s.} Richard Feynman propone la idea teórica de usar
    sistemas cuánticos para simular física cuántica de manera eficiente.
  \item \textbf{1994.} Peter Shor publica un algoritmo cuántico capaz
    de factorizar números enteros exponencialmente más rápido que los
    algoritmos clásicos conocidos. Este resultado amenaza directamente
    a la criptografía RSA actual.
  \item \textbf{1996.} Lov Grover publica un algoritmo cuántico de
    búsqueda con ventaja cuadrática sobre el caso clásico.
  \item \textbf{2019.} Google anuncia haber alcanzado la
    \emph{supremacía cuántica} con su procesador Sycamore de 54
    qubits, aunque el resultado fue disputado por IBM.
\end{itemize}

A pesar del entusiasmo, hay limitaciones importantes que conviene
mencionar. Los qubits son extremadamente frágiles y requieren
temperaturas cercanas al cero absoluto; el fenómeno de la
\textbf{decoherencia} destruye la información cuántica en muy poco
tiempo; y, hasta hoy, todos los dispositivos prácticos pertenecen a la
era \textbf{NISQ} (\emph{Noisy Intermediate-Scale Quantum}), donde los
qubits son pocos y propensos a errores. Las computadoras cuánticas
\emph{no} están destinadas a reemplazar a las computadoras clásicas
para tareas cotidianas; el modelo más probable es que ambas convivan
de forma complementaria.

% ============================================================
\section{Reflexión y conclusiones}
% ============================================================

Resulta fascinante observar cómo la historia de las computadoras refleja un 
proceso de innovación tecnológica que ha transformado radicalmente la sociedad 
en poco más de un siglo. 
A lo largo de las cinco generaciones, se pueden identificar patrones claros: 
cada salto tecnológico ha permitido reducir el costo del cómputo, miniaturizar 
los dispositivos y acercarlos al usuario final, democratizando el acceso a esta 
poderosa herramienta. \\
De igual manera resulta bastante interesante poner en perspectiva el ``de dónde 
venimos'' con el punto en el que nos encontramos actualmente, y el ``hacia dónde 
vamos''.
En ocasiones he escuchado a personas decir: ``la computadora es un dispositivo 
que no es de este planeta'', ``es imposible que los humanos hayamos inventado algo
tan complejo'', y no es para menos, si se piensa en la complejidad de una 
computadora moderna, con sus miles de millones de transistores, su capacidad 
de procesamiento paralelo, y más actualmente, su capacidad de aprendizaje 
automático y de procesamiento de lenguaje natural. Sin embargo, al revisar la 
historia de las computadoras, al intentar entender cómo funcionan (al igual
que todos estos avances tecnológicos), se puede apreciar que no es magia, sino 
el resultado de un proceso acumulativo de innovación, experimentación y 
aprendizaje a lo largo del tiempo por parte del ingenio humano (cientos, miles 
de personas a lo largo de la historia). \\
De esta investigación me llevo una gran expectativa por el futuro de la 
computación, ¿cuál es el siguiente salto tecnológico que nos espera?, y 
¿cómo este va a transformar nuestra sociedad? Me encantaría poder vivirlo y 
aún más, contribuir y ser parte de él.

% ============================================================
\renewcommand{\refname}{Referencias}
\addcontentsline{toc}{section}{Referencias}
% ============================================================

\begin{thebibliography}{9}
\small

\bibitem{tanenbaum2009}
A.~S. Tanenbaum,
\emph{Sistemas Operativos Modernos}, 3.\textsuperscript{a} ed.
México: Pearson Educación, 2009. [Capítulo~1, Sección~1.2: Historia de
los sistemas operativos.]

\bibitem{uvmx}
G.~Hernández,
\emph{Historia de la computación} [en línea].
Universidad Veracruzana, 2011. Disponible en:
\url{https://www.uv.mx/personal/gerhernandez/files/2011/04/historia-compuesta.pdf}

\bibitem{uves}
A.~Fuertes,
\emph{Tema 1: Tecnologías de la información} [en línea].
Universidad de Valencia, Departamento de Informática. Disponible en:
\url{https://www.uv.es/afuertes/Informatica/ficheros/tema1_TI.pdf}

\bibitem{aniei}
ANIEI,
\emph{Curso de Introducción a la Programación: Generaciones de
computadoras} [en línea]. Asociación Nacional de Instituciones de
Educación en Tecnologías de la Información, UAM. Disponible en:
\url{https://aniei.org.mx/paginas/uam/CursoIP/curso_ip_02.html}

\bibitem{ucm}
R.~Carpio López,
\emph{Introducción a la computación cuántica} [trabajo de fin de grado
en línea]. Universidad Complutense de Madrid, 2022. Disponible en:
\url{https://hdl.handle.net/20.500.14352/3166}

\bibitem{scribd}
\emph{Quinta generación de la computadora} [presentación en línea].
Scribd. Disponible en:
\url{https://www.scribd.com/presentation/379887533/Quinta-Generacion-de-La-Computadora}

\end{thebibliography}

\end{document}